% Transcription — fonds Alexandre Grothendieck, shelfmark 18, batch 1 (pages 1-20).
% Established from the facsimile archives/batches/18/batch-01.pdf (a mirror of
% https://grothendieck.umontpellier.fr/18.pdf), per the skill
% transcribe-grothendieck. The batch file carries no cover sheet: PDF page N is
% archivists' page N. This was checked against the pencilled numbers at the
% bottom-left — a "1" on the first leaf and a "20" on the last confirm the
% alignment end to end.
%
% Pass: Opus 5 (claude-opus-5), 2026-09-16 — first pass, unchecked against the pages by a human.
%
% What the batch is. Nine leaves in his hand (pages 4, 5, 8, 10, 11, 12, 14, 16,
% 18, 20), three of them preceded by a cover inscribed with the title of the run
% that follows (pages 3, 9, 13), plus one leaf of a French typescript carrying
% his handwritten corrections (page 6). The remaining leaves — 2, 7, 15, 17, 19,
% 21 — are unrelated typescript (an NSF grant proposal in English, a French
% typescript on formal kernels, a Bourbaki « Tribu » on root systems), several
% of them scanned upside down and none of them marked by him; they are skipped
% without further mention. Page 1 is the same NSF proposal, but he has written
% across the head of it « Papiers donnés à Saav[edra] » — provenance rather than
% mathematics, so the page is skipped too and the inscription recorded only
% here. Nothing on the leaves carries a date and there is no pagination of his.
%
% The mathematics is not one run but three short ones, each opened by its own
% cover, plus one leaf standing on its own. All of it turns on the same question:
% how the several cohomological realisations of a motive — Betti, De Rham,
% Hodge, Tate, crystalline — are to be compared, and what the comparison
% isomorphisms cost in periods.
%
%   3   cover: « Motifs et cristaux sur corps local ».
%   4   two parallel lists of conjectural descriptions of the category of
%       motives: six over a number field (double lattice Betti-De Rham, Betti-
%       Hodge, Hodge structure, Betti-(Gauss)Manin, Betti-Tate, filtered
%       crystalline De Rham), then the primed analogues 3', 2', 5', 6' for
%       p-divisible groups — more generally « pseudo-motifs » — over a local
%       field, with the Hodge-Tate bigrading on N ⊗ C and a filtered crystal
%       with a σ-semi-linear action; closing on « Pourquoi remplacer 1', 4' ?? »
%   5   an extension 1 → G'(Q) → E → π → 1 with two splittings σ₁, σ₂, sitting
%       under a column G(Q_p), G(C), G'(Q), G'(Q_p) and over j₂ : Q_p^* and
%       χ : π; α identifies E with an extension of π by G'(C) and with another
%       by G(C); writing σ̃₁ = σ̃₂ gives σ₁(γ)^{-1}σ₂(γ) = u(γ)·j₂(χ(γ^{-1})),
%       then two remarks on Ker χ and Ker u and the quotient Π₁ = π/Ker u.
%   6   typescript, numbered 18, sections 7.13.10-7.13.12: the specialisation of
%       algebraic cycles is compatible with the ℓ-adic cycle class, but only
%       after Z_ℓ(i) is replaced by Q_ℓ(i); the remark that a torsion-free
%       statement would need a specialisation theory for the Chow ring, which
%       does not exist; the square (7.13.12) for Z^i and Z_ℓ(i). His hand adds
%       « il semble qu'on sera », « des carrés », « lorsque f : X → S est lisse
%       et propre », « le soin », and inks the ℓ's the machine could not type.
%   8   « Twist de Tate sur structures de Hodge mixtes »: the twist E(-1), then
%       E(-n), by W_i(E(-n)) = W_{i+2n}(E) and F^i(E(-n)) = F^{i+n}(E).
%   9   cover: « degrés de transcendance liés aux filtrations et bigraduations
%       (comparaison de Betti, DeRham, Hodge) ».
%  10   the period torsors P = Isom(T_B ⊗ k, T_DR), P' = Isom(T_DR, T_Hdg),
%       P'' = Isom(T_Hdg, T_B ⊗ k), the composition maps P × P' → P'' and its
%       cyclic permutations, α, β, γ with βα = γ^{-1}; the Q-schemes Fil and Big
%       of ⊗-filtrations and ⊗-bigradings of T_B, with Big → Fil a torsor under
%       a unipotent group; a parabolic B' ⊂ G' with unipotent radical U' and
%       B'/U' ≃ L', the central G_m of Tate; an NB proposing the same with F, P
%       in place of Fil, Big and no integrality condition.
%  11   the conjecture deg tr k(g₀) = dim G'' (= dim G − 1) and a corollary; the
%       remark that it contains Schneider's theorem for elliptic curves without
%       complex multiplication, where the transcendence degree in question is 3.
%       The leaf is faint pencil under two long diagonal cancellations and very
%       little of the connective prose carries.
%  12   « Questions sur courbes elliptiques sur Q̄ »: G = Gl(2,Q), or the Weil
%       restriction of G_m from the CM field; g ∈ G(C)/G(Q̄) the Betti-De Rham
%       comparison; three questions a), b), c) — the transcendence degree of
%       Q̄(g), whether the group H generated by the locus T of g is all of G,
%       and whether the period h splitting De Rham is transcendental (it is in
%       Q̄ when C has complex multiplication).
%  13   cover: « Le foncteur de Hodge sur la catégorie des motifs ».
%  14   « Comparaison de classes fondamentales analytiques (= alg.) et
%       topologiques (= transcendantes) »: Y ⊂ X of codimension p, the square
%       comparing H^{2p}_Y(X,Z) → H^{2p}_Y(X,Ω^·_X) with the Poincaré classes
%       P^Z_top, P^C_top, P^Ω_an, and the reduction to Y a point by a transverse
%       section, with the residue df₁∧…∧df_p / f₁…f_p.
%  16   the computation: P_an(1_Y) is the class of dz/z, λ = ∫_C dz/z = 2iπ,
%       hence the boxed P_an((1/2iπ)1_Y) = P_top(1_Y) and its two forms; then
%       the same for H^{i+p}(X) against H^i(Y) with the factor 1/(2iπ)^p.
%  18   over a field k abstractly isomorphic to C: P_alg on a^p(X) is defined
%       without choosing σ, whereas P_σ^{(p)} = (1/2iπ)^p P_alg depends on σ;
%       the case p = 1, X = P^1.
%  20   the same square for i^!(C) → i^!(Ω^·_X); an NB that the analytic direct
%       image does not carry integral classes to integral classes without the
%       factor 1/(2iπ)^p; a boxed square for a^i(X) → H^{2i}(X,Z).
%
% Notation fixed for the batch. Z, Q, C, R, Q̄, Q_p, Z_ℓ, Q_ℓ are set as
% \mathbb{}; k, e and the group letters G, G', G'', H', H'', E, L', U', B' stay
% roman, as do the torsors P, P', P'', Q and the schemes Fil, Big, F. The fibre
% functors are T_{B}, T_{DR}, T_{Hdg}, as he writes them. The De Rham complex is
% \Omega^{\bullet}_{X}. The mixed Hodge structure of page 8 is \mathcal{E}, his
% script E; the group of cycle classes of page 18 and 20 is a^{i}(X), which he
% draws with a rounded, half-calligraphic a. Cycle-class and Poincaré maps keep
% his own labels — P_{\mathrm{top}}, P_{\mathrm{an}}, P_{\mathrm{alg}},
% P^{\mathbb{Z}}_{\mathrm{top}}, P_{\sigma}^{(p)} — and 1_{Y} is the fundamental
% class. He underlines some of his H's once and some twice, evidently to mark
% hypercohomology of a complex against ordinary cohomology, but not consistently
% enough to reconstruct the rule; the transcription writes H throughout and says
% so once, at page 14.
%
% Material facts stated once. Page 5 is written across a landscape leaf, page 11
% is faint pencil under two diagonal cancellations, and pages 10 and 14 are so
% crowded that their marginal annotations are compressed into \note{} rather than
% set out line by line. Page 6 is typed, not handwritten, and is kept because it
% carries his corrections; his insertions are given in the text and identified by
% a \note, the typist's strikings by \struck{}.
%
% Batch boundary. No run crosses into page 21: page 20 closes on a boxed square
% with blank paper below it, and page 21 is unrelated typescript, upside down and
% unmarked.

\documentclass[11pt,a4paper]{article}
% Le préambule commun à toutes les transcriptions du fonds Grothendieck.
%
% Il tient en une idée : **l'incertitude se compose, elle ne se résout pas.**
% Une transcription qui lisse un mot douteux a détruit la seule chose pour
% laquelle on la faisait — un lecteur doit pouvoir savoir, à chaque endroit,
% ce qui était sur le feuillet et ce qui a été ajouté. Les six macros qui
% suivent sont donc l'appareil critique, et non de la décoration.
%
% Les mêmes commandes sont reconnues par `scripts/render.mjs`, qui produit la
% vue de lecture du site. Une macro ajoutée ici doit l'être aussi là-bas,
% faute de quoi le rendu échouera bruyamment — ce qui est voulu : un rendu
% silencieusement faux est pire qu'un rendu absent.

\ProvidesPackage{grothendieck}[2026/08/08 Transcriptions du fonds Grothendieck]

\usepackage{amsmath,amssymb,amsthm}
\usepackage{mathrsfs}
\usepackage{stmaryrd}
% Les diagrammes commutatifs sont partout dans ce fonds : c'est la langue
% même de la géométrie algébrique de Grothendieck, et une transcription qui
% les rendrait en prose ne transcrirait rien.
\usepackage{tikz-cd}
% Français par défaut — c'est la langue des feuillets — mais l'anglais est
% chargé aussi : la traduction et le résumé sont des documents anglais, et la
% typographie française (espaces avant les deux-points) y serait une faute.
% Ces documents basculent par \selectlanguage{english} en tête de corps.
\usepackage[english,french]{babel}
\usepackage{fontspec}
\usepackage{geometry}
\geometry{a4paper,margin=25mm,marginparwidth=28mm}
\usepackage{marginnote}
\usepackage{xcolor}
% ulem, not soul. `soul`'s \st and \ul are built on a character-by-character
% scan that XeLaTeX's font machinery breaks: the document fails with an
% undefined control sequence rather than degrading. ulem does the same two jobs
% and is Unicode-engine safe. `normalem` stops it hijacking \emph, which the
% transcriptions use for Grothendieck's own emphasis.
\usepackage[normalem]{ulem}
\usepackage{hyperref}
\hypersetup{colorlinks=true,linkcolor=black,urlcolor=black,pdfborder={0 0 0}}

% --- Métadonnées du lot -----------------------------------------------------
% Reprises telles quelles par le script de rendu, qui les affiche en tête de
% la vue de lecture. Elles fixent ce que le fichier prétend couvrir : sans
% elles, une transcription flotte sans point d'attache dans un fonds de seize
% mille feuillets.

\newcommand{\folder}[1]{\gdef\grothFolder{#1}}
\newcommand{\batch}[1]{\gdef\grothBatch{#1}}
\newcommand{\leaves}[2]{\gdef\grothFirst{#1}\gdef\grothLast{#2}}
\newcommand{\foldertitle}[1]{\gdef\grothFoldertitle{#1}}
\gdef\grothFolder{?}\gdef\grothBatch{?}\gdef\grothFirst{?}\gdef\grothLast{?}\gdef\grothFoldertitle{}

% --- Le résumé --------------------------------------------------------------
%
% Il ouvre la lecture modernisée, sous le titre « Résumé », et il est composé
% en petit. Ce n'est pas un ornement : c'est le seul passage du document qui
% ne soit pas des mathématiques, et un lecteur doit pouvoir le reconnaître
% avant de l'avoir lu — puis le sauter s'il connaît déjà le sujet, ou s'y
% arrêter s'il ne le connaît pas. Le filet qui le suit marque la reprise du
% corps.

\newenvironment{resume}
  {\small\setlength{\parskip}{0.45em}}
  {\par\medskip\hrule\medskip}

% --- Mots-clés ---------------------------------------------------------------
%
% En anglais, en clôture du résumé de la lecture modernisée : le vocabulaire
% moderne sous lequel chercher ce que les feuillets construisent. Le manifeste
% du site (`npm run manifest`) les extrait pour étiqueter la cote entière —
% cette ligne est la source unique des tags, il n'y a pas de fichier à côté.
\newcommand{\keywords}[1]{\par\smallskip\noindent
  {\footnotesize\textsc{Keywords} --- \textit{#1}}\par}

% --- L'appareil critique ----------------------------------------------------

% Le numéro de feuillet, en marge. C'est l'ancre qui permet de régler un
% désaccord sur une formule en regardant *un* feuillet plutôt qu'une chemise
% de six cents. Le site s'en sert aussi pour tourner les pages du fac-similé
% quand on fait défiler la transcription.
\newcommand{\leaf}[1]{%
  \marginnote{\footnotesize\color{gray!70}\textsf{#1}}%
  \label{leaf:#1}\ignorespaces}

% Illisible : on ne devine pas. Un mot inventé qui se lit comme les autres est
% le pire résultat possible.
\newcommand{\ill}{\textcolor{red!60!black}{[\,\ldots\,]}}

% Lecture douteuse : le mot est donné, et signalé comme incertain.
\newcommand{\uncertain}[1]{\uline{#1}}

% Ajout de l'éditeur — une lettre manquante, un mot sous-entendu.
\newcommand{\add}[1]{\textcolor{blue!50!black}{[#1]}}

% Biffé par Grothendieck lui-même. Ce qu'il a rayé fait partie du document :
% une rature dit souvent où la pensée a changé de direction.
\newcommand{\struck}[1]{\sout{#1}}

% Note du transcripteur, sur l'état matériel du feuillet.
\newcommand{\note}[1]{\par\smallskip{\small\color{gray!60!black}\textit{[#1]}}\par\smallskip}

% Note marginale de Grothendieck, distincte du corps du texte.
\newcommand{\marginal}[1]{\par\smallskip\noindent\hspace{1em}%
  {\small\color{gray!60!black}#1}\par\smallskip}

% --- Titre ------------------------------------------------------------------

\AtBeginDocument{%
  \begin{center}
    {\large\bfseries Fonds Alexandre Grothendieck --- cote n\textsuperscript{o}~\grothFolder}\\[2pt]
    {\itshape\grothFoldertitle}\\[4pt]
    {\small feuillets \grothFirst--\grothLast\ (lot \grothBatch)}\\[2pt]
    {\footnotesize\color{gray!60!black}Transcription établie d'après le fac-similé numérisé
     par l'Université de Montpellier.\\
     Les lectures incertaines sont soulignées, les illisibles notées [\,\ldots\,],
     les ajouts entre crochets.}
  \end{center}
  \vspace{1em}}

\endinput


\folder{18}
\batch{1}
\pages{1}{20}
\dating{[à partir de 1969-vers 1972]}
\watermark{Édition de démonstration}
\foldertitle{Motifs et théorie de Hodge : notes manuscrites (s.d.)}

\begin{document}

\section{Motifs et cristaux sur corps local}

\note{Le titre est celui du feuillet 3, qui ne porte que ces deux lignes,
soulignées, de sa main ; il ouvre la suite des pages 4 à 6. Le feuillet n'est
pas transcrit comme page.}

\page{4}

Descriptions conjecturales des motifs sur corps de nbs \struck{\ill{}} $k$

\begin{enumerate}
  \item Double réseau Betti -- DeRham
  \item Double réseau Betti -- Hodge
  \item \struck{Betti $+$ bigrad.} structure de Hodge \qquad ($k = \overline{\mathbb{Q}}$)
  \item \struck{Betti $+$ \ill{} réseau Betti}
        Betti -- (Gauss)Manin \qquad ($k = \overline{\mathbb{Q}}$)
  \item \struck{Tate} Betti -- Tate
  \item \struck{De \ill{} cristal} De Rham filtré cristallin
\end{enumerate}

Descriptions conjecturales des groupes $p$-divisibles \note{sous « $p$-divisibles »
il insère deux mots dont la lecture reste douteuse, \uncertain{car.\ de $k$}}
(plus gén.\ des pseudo-motifs) \marginal{mod.\ isogénie} sur un corps local
[\struck{\ill{}} d'un \uncertain{anneau} de val.\ discrète complet à corps rés.\
\struck{\ill{}}] :

\begin{enumerate}
  \item[3')] structure de Hodge \struck{(mod.\ ext.\ finie de $K$)}
  \item[2')] double réseau Tate -- Hodge
  \item[5')] Tate (module galoisien $\ell$-adique)
  \item[6')] \struck{module} \struck{Vect.\ de di\ill{}} cristal filtré [vectoriel
        \struck{\ill{}} de dim finie sur $K$, avec une action $\sigma$-semi-linéaire,
        $+$ filtration, graduation]
\end{enumerate}

\marginal{Tate : module \struck{\ill{}} $N$ de t.f.\ sur $\mathbb{Q}_{p}$, et
bigrad.\ sur $N \otimes_{\mathbb{Q}_{p}} C$}

\note{La note marginale ci-dessus est portée au-dessus de l'article 3'). Le
corps des scalaires du produit tensoriel est écrit d'un seul trait et se lit
$\mathbb{Q}_{p}$ plutôt que $\mathbb{Q}$ ; le $C$ qui suit est celui de Tate.}

Pourquoi remplacer 1', 4' ??

\note{Les chiffres 3', 2' et 5' sont écrits par-dessus d'autres, qu'on ne peut
plus lire ; 6' est de premier jet. La liste primée ne comporte donc que quatre
articles sur six, et c'est à quoi répond la question finale.}

\page{5}

\note{Le feuillet est écrit en travers, format à l'italienne. La colonne
verticale de gauche porte, entre $G(\mathbb{Q}_{p})$ et $G(\mathbb{C})$ puis
entre $G(\mathbb{C})$ et $G'(\mathbb{Q})$, deux traits sans pointe dont les
étiquettes se lisent mal : elles sont données ci-dessous comme $p$ et $j_{1}$,
sous toute réserve. Seules les flèches $u$, $j_{2}$ et $\chi$ portent une pointe
nette.}

\begin{tikzcd}[column sep=small, row sep=small, nodes={font=\scriptsize}]
 & \pi \arrow[d, "u"] & & & \\
 & G(\mathbb{Q}_{p}) \arrow[d, no head, "p"'] & & & \\
 & G(\mathbb{C}) \arrow[d, no head, "j_{1}"'] & & & \\
1 \arrow[r] & G'(\mathbb{Q}) \arrow[r] \arrow[d, no head] & E \arrow[r] & \pi \arrow[r] \arrow[l, bend left=30, "\sigma_{1}"] \arrow[l, bend right=30, "\sigma_{2}"'] & 1 \\
 & G'(\mathbb{Q}_{p}) & & & \\
 & G_{m}(\mathbb{Q}_{p}) = \mathbb{Q}_{p}^{*} \arrow[u, "j_{2}"] & & & \\
 & \pi \arrow[u, "\chi"] & & &
\end{tikzcd}

\struck{\ill{}} $\alpha$ permet d'identifier une \uncertain{extension} cent.\ de
$E$ de $\pi$ par $G'(\mathbb{C})$ \ill{} une autre \uncertain{ext.}\ de $\pi$ par
$G(\mathbb{C})$

Le splitting $\sigma_{2}$ de la première \ill{} lui par $\sigma_{2}(R)$ \ill{} :
$G'(\mathbb{Q}_{p})$ et à fortiori : $j_{2}\chi(\pi)$. On trouve donc
\[
  \begin{cases}
    \gamma \mapsto \widetilde{\sigma}_{2}(\gamma) = \sigma_{2}(\gamma) \cdot j_{2}(\chi(\gamma)) \\
    \pi \longrightarrow E
  \end{cases}
\]
un autre splitting.

Le splitting $\sigma_{1}$ dans la seconde \ill{} à $G(\mathbb{Q}_{p})$, et à
fortiori : $u(\gamma)$. On trouve donc
\[
  \gamma \mapsto \sigma_{1}(\gamma) u(\gamma) = \widetilde{\sigma}_{1}(\gamma) .
\]

On écrit que $\widetilde{\sigma}_{1}(\gamma) = \widetilde{\sigma}_{2}(\gamma)$,
i.e.
\[
  \sigma_{2}(\gamma)\, j_{2}(\chi(\gamma)) = \sigma_{1}(\gamma)\, u(\gamma) ,
\]
donc
\[
  \sigma_{1}(\gamma)^{-1} \sigma_{2}(\gamma) = \struck{\ill{}}\; u(\gamma) \cdot j_{2}(\chi(\gamma^{-1})) .
\]

\note{Sous cette dernière ligne il souligne les trois morceaux : le membre de
gauche est noté $\varphi(\gamma)$, et les deux facteurs de droite sont marqués
comme appartenant respectivement à $G(\mathbb{Q}_{p})$ et à
$j_{2}(\mathbb{Q}_{p}^{*})$.}

\begin{enumerate}
  \item[$1^{\circ}$)] Sur $\operatorname{Ker} \chi$, \uncertain{on a}
        $\varphi(\gamma) = u(\gamma)$, donc $u$ \uncertain{commute} sur
        $\operatorname{Ker}$ \struck{\ill{}}
  \item[$2^{\circ}$)] $\operatorname{Ker} u = \struck{\chi(\gamma)} \{\gamma,\ \uncertain{\chi}(\gamma) = 1,\ \varphi(\gamma) = 1\}$
        \uncertain{commun}. Soit $\Pi_{1} = \pi/\operatorname{Ker} u$
        ($\rightarrow$ un gpe analytique $p$-adique, et \ill{}), et
        $\Pi_{1} \xrightarrow{\ \chi_{1}\ } \mathbb{Q}_{p}^{*}$ induit par $p$
\end{enumerate}

\note{Le feuillet s'arrête là : la moitié inférieure est laissée en blanc et
l'argument n'est pas repris sur la page suivante, qui est d'une autre main.}

\page{6}

\note{Ce feuillet est dactylographié — il porte en tête, dans sa propre
pagination, le numéro 18 — et n'est retenu que parce qu'il est corrigé de sa
main. Ses ajouts manuscrits sont donnés entre crochets et signalés ; ce que la
dactylographie a biffé est rendu par \struck{}. La machine n'avait pas de
$\ell$ : les $\ell$ des coefficients sont encrés à la main sur la frappe.}

(7.13.10)

\begin{tikzcd}[column sep=small, row sep=small, nodes={font=\scriptsize}]
\operatorname{Gr}^{i}(X_{\bar{t}})_{\mathbb{Q}} \arrow[r, "\mathrm{cl}"] \arrow[d] & H^{2i}(X_{\bar{t}}, \mathbb{Q}_{\ell}(i)) \arrow[d, "\wr"] \\
\operatorname{Gr}^{i}(X_{\bar{s}})_{\mathbb{Q}} \arrow[r, "\mathrm{cl}"] & H^{2i}(X_{\bar{s}}, \mathbb{Q}_{\ell}(i))
\end{tikzcd}

\struck{7.13.11} Remarque\struck{s} 7.13.11. La commutativité des deux derniers
diagrammes \struck{commutatifs}, qu'on peut exprimer en disant que la
spécialisation des cycles algébriques est compatible avec la formation des
classes de cohomologie $\ell$-adiques associées, n'est pas très satisfaisante
sous la forme énoncée, puisqu'on \struck{\ill{}} a été obligé de remplacer les
coéfficients $\mathbb{Z}_{\ell}(i)$ par $\mathbb{Q}_{\ell}(i)$, c'est à dire de
négliger la torsion. \note{« coéfficients » est de la frappe} Pour avoir un
énoncé qui ne néglige pas la torsion, [il semble qu'on sera] \struck{on sera}
obligé de travailler avec l'anneau de Chow, puisque c'est sur ce dernier que ce
trouve défini naturellement l'homomorphisme $\mathrm{cl}$ comme homomorphisme à
valeurs dans $H^{2*}(-, \mathbb{Z}_{\ell}(*))$, \struck{\ill{}} \struck{\ill{}}
comparer (XIV 4.1). \note{« ce trouve » pour « se trouve », de la frappe}
Malheureusement, on n'a pas développé pour le moment une théorie de la
spécialisation pour l'anneau de Chow \struck{\ill{}} (cf.\ 7.14), qui
permettrait d'écrire [des carrés] \struck{un diagramme} commutatifs analogues à
(7.13.9) et (7.13.10) en termes des anneaux de Chow et de la cohomologie à
coéfficients $\mathbb{Z}_{\ell}(*)$. Tout ce qu'on peut dire, [lorsque
$f \colon X \to S$ est lisse et propre] c'est qu'on a un carré commutatif

(7.13.12)

\begin{tikzcd}[column sep=small, row sep=small, nodes={font=\scriptsize}]
Z^{i}(X_{t}) \arrow[r, "\mathrm{cl}"] \arrow[d, "\sigma"'] & H^{2i}(X_{t}, \mathbb{Z}_{\ell}(i)) \\
Z^{i}(X_{s}) \arrow[r, "\mathrm{cl}"] & H^{2i}(X_{s}, \mathbb{Z}_{\ell}(i)) \arrow[u]
\end{tikzcd}

\note{La flèche verticale de droite est dirigée vers le haut, en sens inverse de
$\sigma$ ; sa pointe est encrée à la main et ne laisse pas de doute.}

\struck{\ill{}} analogue à (7.13.9), et un carré correspondant analogue à
(7.13.10) qu'on laisse au\struck{x} lecteur [le soin] d'écrire, où $Z^{i}$
désigne le groupe des cycles algébriques (modulo rien du tout) de codimension
$i$, \struck{\ill{}} où les flèches horizontales sont les homomorphismes
« classe de cohomologie associée » définie dans (SGA 5 IV), et où $\sigma$
désigne l'homomorphisme de spécialisation des cycles algébriques, associant à un
cycle $Z_{t}$ le cycle $Z_{s}$, où $Z$ est

\note{La phrase s'interrompt au bas du feuillet ; la suite de la dactylographie
ne figure pas dans le dossier.}

\subsection{Twist de Tate sur structures de Hodge mixtes}

\note{Le titre est de lui, porté en tête du feuillet 8 et souligné.}

\page{8}

Soit $\mathcal{E} = (\mathcal{E}, \mathcal{E}_{0}, W_{*}, F^{*})$ une structure
de Hodge mixte. \note{sous la ligne il note que $W_{*}$ est covariante et $F^{*}$
décroissante}

On définit une structure de Hodge $\mathcal{E}(-1)$ (twist de Tate $(1)$), avec
\ill{} $= \mathcal{E}_{0}$ \struck{\ill{}} \marginal{module sur $\mathbb{Z}$},
\ill{} $= \mathcal{E}_{0}$,
\[
  \begin{cases}
    W_{i}(\mathcal{E}(-1)) = W_{i-2}(\mathcal{E}) \\
    F^{i}(\mathcal{E}(-1)) = F^{i-1}(\mathcal{E})
  \end{cases}
\]

\note{Le signe de l'indice de $W$ est repassé et pourrait se lire $+$ ; tel
quel, $W_{i-2}$ s'accorde avec le $F^{i-1}$ de la ligne suivante, mais non avec
la formule générale donnée plus bas pour $\mathcal{E}(-n)$, qui pour $n = 1$
donnerait $W_{i+2}$ et $F^{i+1}$. La discordance est sur le feuillet et n'est
pas corrigée ici.}

(on trouve alors que
$W_{i}(\mathcal{E}(-1)) / W_{i-1}(\mathcal{E}(-1)) \simeq [W_{i}(\mathcal{E}) / W_{i-1}(\mathcal{E})](-1)$)

(twist de Tate sur structures de Hodge \ill{}).

Bien sûr, on définit pour tt entier $n$ un twist $\mathcal{E}(-n)$,
\[
  \begin{cases}
    W_{i}(\mathcal{E}(-n)) = W_{i+2n}(\mathcal{E}) \\
    F^{i}(\mathcal{E}(-n)) = F^{i+n}(\mathcal{E})
  \end{cases}
\]

\note{Le feuillet s'arrête là, la moitié inférieure restant blanche.}

\section{Degrés de transcendance liés aux filtrations et bigraduations (comparaison de Betti, DeRham, Hodge)}

\note{Titre pris au feuillet 9, qui ne porte que ces mots, de sa main, dans
l'angle supérieur droit ; il ouvre la suite des pages 10 à 12. Le feuillet n'est
pas transcrit comme page.}

\page{10}

\note{Le feuillet est écrit sur trois colonnes très serrées, avec des encadrés
et des renvois ; les annotations de marge sont rendues ci-dessous sans chercher
à reproduire leur disposition, et les endroits où le trait ne porte plus sont
laissés \ill{}.}

$G'' \supset H''$ \qquad $\mathbb{Q}$ \qquad $H' \subset G'$ \uncertain{$\supset$} $G_{k}$

\[
  \begin{cases}
    P = \operatorname{Isom}(T_{B} \otimes k,\ T_{DR}) \\
    Q \subset P' = \operatorname{Isom}(T_{DR},\ T_{Hdg}) \simeq Q \times^{H'} G' \\
    P'' = \operatorname{Isom}(T_{Hdg},\ T_{B} \otimes k)
  \end{cases}
  \qquad \uncertain{\text{déf.}} / k
\]

\[
  \begin{cases}
    \mathrm{Fil} = \operatorname{Filt}_{\otimes,\ \ill{}}(T_{B}) \\
    \mathrm{Big} = \operatorname{Bigr}_{\otimes,\ \ill{}}(T_{B})
  \end{cases}
  \qquad \uncertain{\text{déf.}} / \mathbb{Q}
\]

\note{Les indices des deux foncteurs $\operatorname{Filt}$ et
$\operatorname{Bigr}$ portent une condition d'appartenance à une classe liée à
la filtration, resp.\ à la bigraduation, de $T_{DR}$, que l'écriture ne permet
pas de restituer.}

\[
  P \times P' \to P'', \qquad P' \times P'' \to P, \qquad P'' \times P \to P',
  \qquad Q \subset P' .
\]

$\mathrm{Big} \xrightarrow{\ p\ } \mathrm{Fil}$ fibration, faisant de
$\mathrm{Big}$ un torseur sur $\mathrm{Fil}$, de gr.\ \ill{}
$U_{\mathrm{Fil}}$, un gpe unipotent \ill{}

$\alpha \in P(\mathbb{C})$, $\beta \in Q(\mathbb{C}) \subset P'(\mathbb{C})$,
$\gamma \in P''(\mathbb{C})$ :

\begin{tikzcd}[column sep=small, row sep=small, nodes={font=\scriptsize}]
T_{B} \otimes_{\mathbb{Q}} \mathbb{C} \arrow[r, "\alpha"] & T_{DR} \otimes_{k} \mathbb{C} \arrow[d, "\beta"] \\
 & T_{Hdg} \otimes_{\mathbb{Q}} \mathbb{C} \arrow[ul, "\gamma"]
\end{tikzcd}

\[
  \beta\alpha = \gamma^{-1}, \qquad \gamma\beta = \alpha^{-1}, \qquad \alpha\gamma = \beta^{-1} .
\]

$U' \subset B' \subset G'$ ; \uncertain{rad.\ unipotent} associé à la filtration
de $T_{DR}$ ; $Q = $ \ill{} des sous-groupes de Lévi de $B'$ ;
$B'/U' \simeq L' \longleftarrow G_{m}$ (hom.\ central) (Tate).

N.B. Soient $F \supset \mathrm{Fil}$, $P \supset \mathrm{Big}$ définis comme
ci-contre, mais sans condition d'\uncertain{intégrance}. Alors \ill{} et des
schémas \struck{\ill{}} $A = F/G$, $B = P/G$ \ill{} sur $\mathbb{Q}$. Alors
$\varphi \in F(\mathbb{C})$, $\psi \in P(\mathbb{C})$ \ill{}

\page{11}

\note{Feuillet au crayon, très pâle, barré de deux longues diagonales. Les
énoncés portent, la prose de liaison presque pas : ce qui suit est ce que le
trait soutient, le reste est laissé \ill{}.}

\ill{} mettons \uncertain{$(1, i, j)$} sont déterminés, je crois. Reste \ill{}
pas. Nous allons examiner le degré de transcendance.

Conjecture : si $k$ \ill{}, alors
\[
  \operatorname{deg\,tr} k(g_{0}) = \dim G'' \quad (= \dim G - 1) .
\]

\ill{} $g_{0}$ est un point \uncertain{générique} de $G$ \ill{}

Corollaire. $G''$ est le \uncertain{groupe algébrique} \ill{}

\ill{} contient le th.\ de Schneider, pour les courbes elliptiques sans mult.\
complexe. En fait, il signifie que le deg.\ tr.\ \ill{} en question est $3$ (les
seules relations étant \ill{}).

\note{Une dernière ligne, en partie couverte par les diagonales, énumère des
quantités du type $\zeta^{-1}\omega_{1}$, $\zeta^{-1}\omega_{2}$, dont ni les
indices ni les exposants ne se lisent.}

\subsection{Questions sur courbes elliptiques sur $\overline{\mathbb{Q}}$}

\note{Le titre est de lui, porté en tête du feuillet 12 et souligné.}

\page{12}

Soit $C$ une telle courbe, $G$ son gpe de Galois \uncertain{motivique}
[$G = Gl(2, \mathbb{Q})$ ou $G = \prod_{K/\mathbb{Q}} G_{m}$, $K$ le corps de
multiplication complexe]. Soit $g \in G(\mathbb{C})/G(\overline{\mathbb{Q}})$
définissant \ill{} des \uncertain{réseaux} transcendant de Betti \ill{} celui de
De Rham.

\ill{} degré de transcendance égal à $\dim G$ ($4$, ou $2$) ?

\begin{enumerate}
  \item[a)] $\overline{\mathbb{Q}}(g)$ a degré de transcendance \ill{}
  \item[b)] Soit $T$ le lieu de $g$ sur $G_{\overline{\mathbb{Q}}}$,
        \struck{\ill{}} \ill{} \uncertain{contenant} $e$, et soit $H$ le gpe
        engendré par $T$. \ill{} A-t-on \struck{\ill{}} $H = G$ \ill{} ?
  \item[c)] Soit \struck{$g_{R} H \subset G'_{\mathbb{Q}}$ \ill{}} \ill{}
        (\uncertain{additif}) définissant le splitting transcendant \ill{} de De
        Rham. Peut-on avoir $h \in \overline{\mathbb{Q}}$, ou $h$ est-il toujours
        transcendant sur $\mathbb{Q}$ ? [Cas où $C$ n'a pas de
        multiplication complexe]
\end{enumerate}

\marginal{si $C$ a de la multiplication complexe, on a $h \in \overline{\mathbb{Q}}$ !}

\note{Cette dernière note est entourée d'un trait, dans la marge gauche, en face
de la question c). La ligne « A-t-on $H = G$ » est elle-même rayée d'un long
trait, de même que le début de c) ; la question subsiste néanmoins, puisqu'il la
réécrit en dessous.}

\section{Le foncteur de Hodge sur la catégorie des motifs}

\note{Titre pris au feuillet 13, de sa main ; il se poursuit par
« (Comparaison des deux \uncertain{fond.\ alg.}\ et transcendants) », dont les
deux mots abrégés ne se lisent pas avec certitude. Le feuillet ouvre la suite
des pages 14 à 20 et n'est pas transcrit comme page.}

\subsection{Comparaison de classes fondamentales analytiques (= alg.) et topologiques (= transcendantes)}

\note{Le titre est de lui, en tête du feuillet 14 et souligné.}

\page{14}

$X$ variété \struck{analyti\ill{}} holomorphe, $Y$ ss-variété holomorphe de
codim $p$ partout, $\Omega^{\bullet}_{X}$ résolution de $\mathbb{C}$.

\note{Il souligne d'un trait, et parfois de deux, certains des $H$ qui suivent,
apparemment pour distinguer l'hypercohomologie d'un complexe de la cohomologie
ordinaire ; l'usage n'est pas assez régulier pour qu'on puisse en tirer une
règle, et la transcription écrit $H$ partout.}

\begin{tikzcd}[column sep=small, row sep=small, nodes={font=\scriptsize}]
H^{2p}_{Y}(X, \mathbb{Z}) \arrow[r, "i"] & H^{2p}_{Y}(X, \mathbb{Z}) \otimes_{\mathbb{Z}} \mathbb{C} \arrow[r, "\mathrm{can}"] & H^{2p}_{Y}(X, \mathbb{C}) \arrow[r, "\sim"] & H^{2p}_{Y}(X, \Omega^{\bullet}_{X}) \\
\mathbb{Z}_{Y} \arrow[u, "P^{\mathbb{Z}}_{\mathrm{top}}"] \arrow[r, hook] & \mathbb{C}_{Y} & \mathbb{C}_{Y} \arrow[u, "P^{\mathbb{C}}_{\mathrm{top}}"] & \mathbb{C}_{Y} \arrow[u, "P^{\Omega}_{\mathrm{an}}"]
\end{tikzcd}

\note{Le diagramme est resserré et son rang inférieur n'est pas aligné colonne
par colonne sur le rang supérieur : il est ici régularisé. La flèche du haut à
droite est étiquetée « De Rham (i.e.\ can) » ; celle du bas à gauche
« topologique / can ». Les trois flèches verticales portent, outre leur nom, une
marque d'isomorphisme. À droite du diagramme il note $H^{p}_{Y}(X, \Omega^{p}_{X})$
et le mot « résidu », et une flèche de facteur $1/(2i\pi)^{p}$ ; un
$H^{p}_{Y}(X, \Omega^{\bullet}_{X})$ du premier jet est raturé au milieu.}

On est ramené au cas où $Y$ est réduit à un pt.

(prendre une \struck{section} $X'^{n-p}$ transversale : $Y$ en est un pt) et par
\uncertain{comp.-produit} au cas \ill{}

\marginal{Poincaré analytique [dépend d'un \uncertain{ordre} en alg.]}

\note{Dans la marge droite il écrit que si $f_{1}, \ldots, f_{p}$ est un système
d'équations locales de $Y$, on prend la section
$df_{1} \wedge \cdots \wedge df_{p} / f_{1} \cdots f_{p}$ de
$H^{p}_{Y}(\Omega^{p}_{X})$ ; plusieurs mots de ce passage restent douteux.}

\page{16}

Alors on voit que $P_{\mathrm{an}}(1_{Y})$ est la classe dans le \ill{} sur
$X - a$ représentée par la forme diff.\ $dz/z$, et par suite \ill{} si
$P_{\mathrm{an}}(1_{Y}) = \lambda P_{\mathrm{top}}(1_{Y})$, \ill{}
\[
  \lambda = \int_{C} P_{\mathrm{an}}(1_{Y}) = \int_{C} \frac{dz}{z} = 2i\pi .
\]

\note{Deux flèches montent vers le $C$ des deux intégrales : la première le
qualifie de cercle orienté, la seconde renvoie à Betti ; le mot qui qualifie
l'orientation ne se lit pas avec certitude.}

Donc
\[
  \boxed{\ P_{\mathrm{an}}\!\left(\tfrac{1}{2i\pi}\, 1_{Y}\right) = P_{\mathrm{top}}(1_{Y})\ }
\]

\note{Il encadre cette formule, puis la réécrit sous ses deux formes dans un
second encadré.}

\[
  \begin{cases}
    P_{\mathrm{top}}(1_{Y}) = \dfrac{1}{2i\pi}\, P_{\mathrm{an}}(1_{Y}) \\[4pt]
    P_{\mathrm{an}}(1_{Y}) = 2i\pi\, P_{\mathrm{top}}(1_{Y})
  \end{cases}
\]

Ex. Considérons le diagramme

\begin{tikzcd}[column sep=small, row sep=small, nodes={font=\scriptsize}]
H^{i+p}(X, \mathbb{Z}) \arrow[r] & H^{i+p}(X, \mathbb{C}) \arrow[r, "\sim"] & H^{i+p}(X, \Omega^{\bullet}_{X}) \\
H^{i}(Y, \mathbb{Z}) \arrow[u, "P_{\mathrm{top}}"] \arrow[r] & H^{i}(Y, \mathbb{C}) \arrow[u, "P_{\mathrm{top}}"] \arrow[r, "\frac{1}{(2i\pi)^{p}}\,\mathrm{can}"] & H^{i}(Y, \Omega^{\bullet}_{Y}) \arrow[u, "P_{\mathrm{an}}"]
\end{tikzcd}

\note{Les indices des $H$ du rang inférieur sont surchargés et illisibles ; ils
portent vraisemblablement un support. Au-dessus de la flèche du bas à droite un
mot du premier jet est biffé.}

\page{18}

Soit $X$ un schéma algébrique lisse sur un corps $k$ isomorphe à $\mathbb{C}$
[de façon non \uncertain{primitive}]. Alors
\[
  a^{p}(X) \xrightarrow{\ P_{\mathrm{alg}}\ } H^{2p}(X, \Omega^{\bullet}_{X})
\]
est bien définie. Le choix d'un isom.\ $k \xrightarrow{\ \sim\ } \mathbb{C}$
\uncertain{permet} en outre de définir
\[
  P_{\sigma}^{(p)} \colon a^{p}(X) \xrightarrow{\ \mathrm{can} \circ P_{\mathrm{top}}\ } H^{2p}(X, \Omega^{\bullet}_{X})
\]
qui se factorise par un \uncertain{morphisme} \ill{}
$H^{2p}(X_{\sigma}, \mathbb{Z}) \to H^{2p}(X, \Omega^{\bullet}_{X})$. \ill{}

\marginal{can $=$ De Rham $+$ \ill{}}

\[
  P_{\sigma}^{(p)} = \left(\frac{1}{2i\pi}\right)^{p} P_{\mathrm{alg}} ,
\]
ce qui montre que \ill{}

$P_{\sigma}^{(p)}$ dépend beaucoup de $\sigma$. (Dans le cas où
$a^{p}(X) \xrightarrow{\ P_{\sigma}^{(p)}\ } H^{2p}(X_{\sigma}, \mathbb{Z})$ est
un isom.\ (ex.\ : $p = 1$, $X = \mathbb{P}^{1}$) \ill{})

\note{Le bas du feuillet reprend le même calcul avec un facteur
$1/(2i\pi)^{?}$ appliqué à $P_{\sigma}(a^{p}(X))$ et une quantité notée
$E_{\sigma}^{2i}$ ; ni l'exposant ni la définition de $E_{\sigma}$ ne se lisent.
Une note de marge porte, soulignés, les mots « De Rham $+$ \ill{} » et
« $\pi$ \ill{} transcendantal ».}

\page{20}

\begin{tikzcd}[column sep=small, row sep=small, nodes={font=\scriptsize}]
H^{*}(X, \mathbb{C}) \arrow[r, "\sim"] & H^{*}(X, \Omega^{\bullet}_{X}) \\
{} \arrow[u, "\wr"] & {} \arrow[u]
\end{tikzcd}

\note{Le rang inférieur de ce premier diagramme est couvert par le diagramme
suivant et ne se lit pas ; la flèche horizontale du bas porte l'étiquette
$\omega$ et la mention « (trivial) ».}

\begin{tikzcd}[column sep=small, row sep=small, nodes={font=\scriptsize}]
H^{*}(Y, i^{!}(\mathbb{C})) \arrow[r, "\sim"] & H^{*}(Y, i^{!}(\Omega^{\bullet}_{X})) \\
H^{*}(Y, \mathbb{C}) \arrow[u, "P_{\mathrm{top}}"] \arrow[r, "\frac{1}{2i\pi}\,\mathrm{can}"] & H^{*}(Y, \Omega^{\bullet}_{Y}) \arrow[u, "P_{\mathrm{an}}"]
\end{tikzcd}

qui provient de la \ill{} can.\ \ill{}

\begin{tikzcd}[column sep=small, row sep=small, nodes={font=\scriptsize}]
i^{!}(\mathbb{C}) \arrow[r, "\sim"] & i^{!}(\Omega^{\bullet}_{X}) \\
\mathbb{C}_{Y} \arrow[u, "P_{\mathrm{top}}"] \arrow[r, "\frac{1}{(2i\pi)^{p}}\,\mathrm{can}"] & \Omega^{\bullet}_{Y} \arrow[u]
\end{tikzcd}

lui \ill{} contenu dans la formule fondamentale \ill{}

N.B. On voit que l'image directe analytique ne transforme pas classes entières
en classes entières : il faut un facteur $\dfrac{1}{(2i\pi)^{p}}$ pour qu'il en
soit \uncertain{ainsi}.

\begin{tikzcd}[column sep=small, row sep=small, nodes={font=\scriptsize}]
a^{i}(X) \arrow[r, "P^{\mathbb{Z}}_{\mathrm{top}}"] \arrow[d, "(\frac{1}{2i\pi})^{p} P_{\mathrm{an}}"'] & H^{2i}(X, \mathbb{Z}) \arrow[d, "\mathrm{can}"] \\
H^{2i}(X, \Omega^{\bullet}_{X}) \arrow[r, "\sim"] & H^{2i}(X, \mathbb{C})
\end{tikzcd}

\note{Ce dernier diagramme est encadré. À gauche il rappelle que $a^{i}(X)$ est
le groupe des cycles analytiques de codimension $i$, et à droite il note que le
carré est commutatif. L'exposant du facteur $1/(2i\pi)$ est écrit $p$ alors que
le reste du diagramme est indexé par $i$ ; la discordance est sur le feuillet.}

\note{Le feuillet s'arrête sur cet encadré, et la page suivante n'est pas de sa
main : l'argument ne se poursuit pas dans ce lot.}

\end{document}
